\part{Strategic Behaviour and Extensions}

\chapter{Game Theory I: Strategic-Form Games}
\label{chap:game-theory-static}

\section{Learning objectives}

After completing this chapter, you should be able to identify players, strategies, and payoffs; find dominant strategies and Nash equilibria; and evaluate equilibrium efficiency.

\section{Games and strategy}

\keyterm{Game theory} studies situations in which each person's outcome depends on their own action and the actions of others. A game specifies players, available strategies, payoffs, and rules. A \keyterm{strategic-form game} represents these elements in a payoff matrix.

A strategy is \keyterm{dominant} if it gives a player a higher payoff than every alternative regardless of what the other player does. A dominated strategy is never a best response and can often be eliminated from analysis.

\section{Nash equilibrium}

A \keyterm{Nash equilibrium} is a profile of strategies in which every player's choice is a best response to the choices of others. No player can gain by changing their action unilaterally. Nash equilibrium predicts stability, not necessarily cooperation, fairness, or efficiency.

The Prisoner's Dilemma illustrates this distinction. Each player has an incentive to choose a self-protective action, even though both would be better off if they cooperated. The model applies to price competition, arms races, environmental agreements, and many political interactions.

\section{Pareto efficiency}

An outcome is \keyterm{Pareto efficient} if no other feasible outcome makes someone better off without harming anyone else. A Nash equilibrium can be Pareto inefficient, as in the Prisoner's Dilemma. This explains why institutions, contracts, repeated interaction, and norms can be valuable: they may alter incentives so that individually rational choices are more socially desirable.

\section*{Chapter Review}

\subsection*{Key concepts}
\begin{itemize}
  \item \textbf{Player}: decision-maker in a game.
  \item \textbf{Strategy}: complete action available to a player.
  \item \textbf{Dominant strategy}: best regardless of opponents' actions.
  \item \textbf{Best response}: optimal response to a given opponent action.
  \item \textbf{Nash equilibrium}: mutual best responses.
\end{itemize}

\section*{Exercises}

\subsection*{Conceptual questions}
\begin{enumerate}[label=\arabic*.]
  \item Define a game and a Nash equilibrium.
  \item Why need a Nash equilibrium not be Pareto efficient?
  \item Distinguish a dominant strategy from a best response.
\end{enumerate}

\subsection*{Graphical exercises}
\begin{enumerate}[resume,label=\arabic*.]
  \item Construct a $2\times2$ payoff matrix for a Prisoner's Dilemma and mark the Nash equilibrium.
\end{enumerate}

\subsection*{Analytical exercises}
\begin{enumerate}[resume,label=\arabic*.]
  \item In the matrix with $(A1,B1)=(4,4)$, $(A1,B2)=(1,5)$, $(A2,B1)=(5,1)$, and $(A2,B2)=(2,2)$, identify dominant strategies, Nash equilibrium, and Pareto-efficient outcomes.
  \item Give a political or economic example in which individually rational choices create an inefficient collective result.
\end{enumerate}

\subsection*{Exam-style multiple-choice questions}
\begin{enumerate}[resume,label=\arabic*.]
  \item A Nash equilibrium is an outcome where: (a) every player cooperates; (b) no player benefits by unilaterally changing action; (c) all payoffs are equal; or (d) one player wins.
  \item A dominant strategy is: (a) optimal only if others cooperate; (b) better regardless of opponents' actions; (c) always Pareto efficient; or (d) unavailable in $2\times2$ games.
\end{enumerate}

\subsection*{PPE discussion question}

\noindent When a Nash equilibrium is socially undesirable, should policy change payoffs, restrict strategies, or encourage cooperation? Discuss using one concrete example.
